\documentclass[11pt]{article}
\usepackage[a4paper,margin=22mm]{geometry}
\usepackage{fontspec}
\setmainfont{DejaVu Serif}
\setsansfont{DejaVu Sans}
\usepackage{microtype}
\usepackage{xcolor}
\usepackage{hyperref}
\usepackage{enumitem}
\usepackage{parskip}
\definecolor{Accent}{HTML}{971D32}
\definecolor{Muted}{HTML}{666666}
\hypersetup{colorlinks=true,urlcolor=Accent,linkcolor=Accent}
\setlist{leftmargin=1.6em,itemsep=0.3em,topsep=0.35em}
\setcounter{tocdepth}{2}
\renewcommand{\contentsname}{Contents}
\newcommand{\sectionrule}{\par\vspace{-0.5em}\noindent\textcolor{Accent}{\rule{\linewidth}{0.6pt}}\par}
\begin{document}

\begin{center}
{\sffamily\bfseries\color{Accent} TOEFL WORD OF THE DAY\par}
\vspace{8mm}
{\fontsize{42}{48}\selectfont\bfseries cogent\par}
\vspace{2mm}
{\Large\sffamily\color{Accent} /ˈkoʊdʒənt/\par}
\vspace{2mm}
{\small\sffamily Local audio phrase: cogent\par}
{\small\sffamily Audio file: audios/cogent.mp3\par}
\vspace{2mm}
{\large\sffamily\color{Muted} adjective\par}
\vspace{5mm}
{\sffamily clear, logical, and convincing\par}
\vspace{6mm}
{\large Cogent describes reasoning that is expressed clearly and is strong enough to persuade people.\par}
\vspace{6mm}
\begin{minipage}{0.84\textwidth}
\itshape “The researcher presented a cogent argument for expanding access to public transportation.”
\end{minipage}\par
\vspace{10mm}
\textbf{Date:} September 14th 2026\quad\textbullet\quad \textbf{Level:} B1--mid-B2 TOEFL
\vfill
Created and compiled by \textbf{Amir Kooshky}\par
\vspace{2mm}
\href{https://t.me/KooshkyTOEFL}{Telegram Channel}\quad\textbullet\quad
\href{https://www.instagram.com/kooshkytoefl}{Instagram}\quad\textbullet\quad
\href{https://t.me/KooshkyTOEFL\_pv}{Personal Telegram}
\end{center}
\newpage
\tableofcontents
\newpage

\section{Meanings and usage}
\sectionrule
\subsection{Meaning 1: Clear, logical, and convincing}
\textbf{Part of speech:} adjective

\textbf{IPA:} /ˈkoʊdʒənt/

\textbf{Pronunciation phrase:} cogent

\textbf{Local audio file:} audios/cogent.mp3

\textbf{Register:} formal; common in academic, professional, legal, and critical discussion

\textbf{Definition:} clearly expressed, logically strong, and likely to persuade people that an idea is true or reasonable

\textbf{In simpler words:} clear and convincing because it makes good sense

\textbf{Typical pattern:} a cogent + argument, reason, explanation, analysis, or case

\subsubsection{Natural examples}
\begin{enumerate}
  \item The researcher presented a cogent argument for expanding access to public transportation.
  \item Her essay offers a cogent explanation of why the policy produced unexpected results.
  \item The committee found the economic evidence cogent enough to justify further investment.
  \item He gave several cogent reasons for preserving the historic building.
  \item A cogent analysis must connect each conclusion to reliable evidence.
  \item Although the proposal was controversial, its supporters made a cogent case for reform.
  \item The speaker's response was brief but cogent, addressing the central objection directly.
  \item Without accurate data, even a confidently delivered argument may not be cogent.
\end{enumerate}

\subsubsection{Nuance and implication}
\textbf{Central nuance:} Cogent combines clarity with persuasive logical force. Something is not cogent merely because it is easy to understand.

\textbf{Usual implication:} The reasoning is organized well and supported strongly enough to influence a thoughtful reader or listener.

\textbf{Compared with coherent:} A coherent argument is orderly and internally connected, but it may still rely on weak evidence. A cogent argument is coherent and also convincing.

\subsubsection{Grammar notes}
\textbf{How to use it:} Use it before a noun, as in a cogent argument, or after a linking verb, as in Her explanation was cogent.

\textbf{What can follow it:} It most often describes reasoning or communication, including an argument, reason, case, explanation, analysis, response, or critique.

\textbf{Position in a sentence:} Words such as highly, particularly, especially, and remarkably can appear before cogent.

\subsubsection{Useful patterns}
\textbf{Pattern:} a cogent argument for or against + noun or -ing form

\textbf{Use:} Use this to describe persuasive reasoning supporting or opposing a position.

\textbf{Example:} The report makes a cogent argument for investing in preventive health care.

\textbf{Pattern:} a cogent reason for + noun or -ing form

\textbf{Use:} Use this for a clearly explained reason that provides strong justification.

\textbf{Example:} The authors give no cogent reason for excluding rural communities from the study.

\textbf{Pattern:} make or present a cogent case

\textbf{Use:} Use this when several connected reasons support a conclusion or proposal.

\textbf{Example:} The evidence presents a cogent case for revising the current guidelines.

\textbf{Pattern:} a cogent explanation or analysis

\textbf{Use:} Use this for an explanation or analysis that is clear, logical, and persuasive.

\textbf{Example:} She offered a cogent analysis of the relationship between housing costs and migration.

\subsubsection{Common wrong usage}
\paragraph{Correction 1}
\textbf{Wrong:} The instructions were cogent, so everyone could follow the steps.

\textbf{Correct:} The instructions were clear, so everyone could follow the steps.

\textbf{Why:} Use cogent when ideas or reasons are persuasive as well as clear. For ordinary instructions that are easy to understand, clear is more natural.

\paragraph{Correction 2}
\textbf{Wrong:} The researcher cogently argumented for a larger sample.

\textbf{Correct:} The researcher argued cogently for a larger sample.

\textbf{Why:} Argue is the verb. Cogently is an adverb that can describe how someone argues.

\paragraph{Correction 3}
\textbf{Wrong:} She provided a cogent evidence for her conclusion.

\textbf{Correct:} She provided cogent evidence for her conclusion.

\textbf{Why:} Evidence is normally uncountable, so do not use a before it.

\section{Word family}
\sectionrule
\subsection{cogency}
\textbf{Part of speech:} uncountable noun

\textbf{IPA:} /ˈkoʊdʒənsi/

\textbf{Definition:} the quality of being clear, logical, and convincing

\textbf{Relationship to the headword:} This noun names the convincing logical strength described by cogent.

\textbf{Grammar pattern:} the cogency of + argument, explanation, or analysis

\textbf{Usage note:} It is formal and normally used without a or an.

\subsubsection{Examples}
\begin{enumerate}
  \item The cogency of her argument depended on the strength of the supporting evidence.
  \item Reviewers praised the report for its clarity and cogency.
\end{enumerate}

\subsection{cogently}
\textbf{Part of speech:} adverb

\textbf{IPA:} /ˈkoʊdʒəntli/

\textbf{Definition:} in a clear, logical, and convincing way

\textbf{Relationship to the headword:} This adverb describes presenting ideas in a cogent way.

\textbf{Grammar pattern:} cogently + argue, explain, demonstrate, or present

\textbf{Usage note:} It is formal and most often modifies verbs connected with reasoning or explanation.

\subsubsection{Examples}
\begin{enumerate}
  \item The author cogently explains why the two variables should be studied separately.
  \item She argued cogently for a more flexible approach.
\end{enumerate}

\section{Common collocations}
\sectionrule
\subsection{cogent argument}
\textbf{Applies to:} Clear, logical, and convincing

\textbf{Meaning:} an argument that is clear, logical, and persuasive

\textbf{Example:} The article develops a cogent argument against relying on a single measure of intelligence.

\subsection{cogent reason}
\textbf{Applies to:} Clear, logical, and convincing

\textbf{Meaning:} a reason that provides strong and persuasive justification

\textbf{Example:} There is a cogent reason for collecting data over a longer period.

\subsection{cogent explanation}
\textbf{Applies to:} Clear, logical, and convincing

\textbf{Meaning:} an explanation that is easy to follow and supported by sound reasoning

\textbf{Example:} The theory offers a cogent explanation for the observed pattern.

\subsection{cogent case}
\textbf{Applies to:} Clear, logical, and convincing

\textbf{Meaning:} a persuasive set of reasons supporting a position

\textbf{Pattern:} make or present a cogent case for or against + noun or -ing form

\textbf{Example:} The researchers make a cogent case for changing the way the disease is monitored.

\subsection{cogent analysis}
\textbf{Applies to:} Clear, logical, and convincing

\textbf{Meaning:} a clear and well-reasoned examination of a subject

\textbf{Example:} Her cogent analysis revealed several weaknesses in the original study.

\section{Synonyms and antonyms}
\sectionrule
\subsection{Close synonyms}
\subsubsection{convincing}
\textbf{Part of speech:} adjective

\textbf{Applies to:} Clear, logical, and convincing

\textbf{Shared meaning:} Both can describe something persuasive enough to produce belief.

\textbf{Difference:} Convincing is broader and can describe anything that causes belief, including a performance or imitation. Cogent is more formal and usually describes reasoning or its expression.

\textbf{Example:} The researchers provided convincing evidence that the intervention was effective.

\subsubsection{compelling}
\textbf{Part of speech:} adjective

\textbf{Applies to:} Clear, logical, and convincing

\textbf{Shared meaning:} Both can describe reasons or evidence that are difficult to dismiss.

\textbf{Difference:} Compelling emphasizes powerful persuasive force and can also describe something that strongly attracts attention. Cogent emphasizes clear, sound reasoning.

\textbf{Example:} The documentary presents compelling evidence of environmental change.

\subsubsection{persuasive}
\textbf{Part of speech:} adjective

\textbf{Applies to:} Clear, logical, and convincing

\textbf{Shared meaning:} Both describe communication that can lead people to accept a position.

\textbf{Difference:} Persuasive focuses on the ability to influence someone. Cogent additionally suggests that the persuasion comes from clarity and logical strength.

\textbf{Example:} Her presentation was persuasive, and the committee approved the proposal.

\subsection{Direct or useful antonyms}
\subsubsection{unconvincing}
\textbf{Part of speech:} adjective

\textbf{Applies to:} Clear, logical, and convincing

\textbf{Opposition:} An unconvincing argument fails to persuade, while a cogent argument succeeds through clear and strong reasoning.

\textbf{Example:} The explanation was unconvincing because it ignored several important facts.

\subsubsection{incoherent}
\textbf{Part of speech:} adjective

\textbf{Applies to:} Clear, logical, and convincing

\textbf{Opposition:} Incoherent reasoning is unclear or poorly connected, while cogent reasoning is clearly organized and logically persuasive.

\textbf{Usage note:} Incoherent is a direct opposite mainly when cogent refers to clarity and logical organization.

\textbf{Example:} The report's conclusion was incoherent and did not follow from the evidence.

\section{Words learners may confuse}
\sectionrule
\subsection{coherent}
\textbf{Part of speech:} adjective

\textbf{Why learners confuse them:} Both words can praise the clarity and organization of ideas.

\textbf{Key difference:} Coherent means clear and logically connected. Cogent goes further: the reasoning is also strong enough to convince.

\textbf{cogent:} The lawyer made a cogent argument supported by reliable evidence.

\textbf{coherent:} The essay was coherent, but its main claim lacked supporting evidence.

\subsection{cognitive}
\textbf{Part of speech:} adjective

\textbf{Why learners confuse them:} The words look somewhat similar but have different meanings and origins.

\textbf{Key difference:} Cogent describes clear and convincing reasoning. Cognitive relates to mental processes such as thinking, learning, memory, and perception.

\textbf{cogent:} The report offers a cogent explanation of the results.

\textbf{cognitive:} Sleep can affect cognitive performance and memory.

\section{Etymology}
\sectionrule
\textbf{Origin:} Cogent comes through French from the Latin verb cogere, meaning to drive or force together.

\textbf{Development:} The idea of force developed into the modern sense of reasoning that strongly moves someone toward a conclusion.

\textbf{Memory link:} A cogent argument drives its reasons together into one convincing conclusion.

\section{Practice exercises}
\sectionrule
\subsection{Word family choice}
Choose the best member of the word family for each blank.

\begin{enumerate}
  \item The author \_\_\_\_\_ explains how the new evidence changes the debate.
  \item The \_\_\_\_\_ of the argument depends on the quality of its evidence.
  \item The report presents a \_\_\_\_\_ case for reducing unnecessary energy use.
\end{enumerate}

\subsection{Correct or incorrect?}
Decide whether each sentence is correct. Correct the inaccurate ones.

\begin{enumerate}
  \item The researcher offered a cogent explanation supported by several independent studies.
  \item The map was cogent, so tourists could easily find the station.
  \item The committee found her reasons highly cogent.
  \item He gave a cogently argument for changing the law.
\end{enumerate}

\subsection{Collocation practice}
Complete each sentence with a natural collocation.

\begin{enumerate}
  \item The article makes a cogent \_\_\_\_\_ for protecting the wetland.
  \item The researchers offered a cogent \_\_\_\_\_ of the unexpected results.
  \item The report presents a cogent \_\_\_\_\_ for increasing the sample size.
\end{enumerate}

\subsection{Complete answer key}
\subsubsection{Word family choice}
\begin{enumerate}
  \item \textbf{cogently.} The author cogently explains how the new evidence changes the debate. \emph{Use the adverb cogently to describe how the author explains.}
  \item \textbf{cogency.} The cogency of the argument depends on the quality of its evidence. \emph{Use the noun cogency after the.}
  \item \textbf{cogent.} The report presents a cogent case for reducing unnecessary energy use. \emph{Use the adjective cogent before the noun case.}
\end{enumerate}

\subsubsection{Correct or incorrect?}
\begin{enumerate}
  \item \textbf{Correct} \emph{Cogent naturally describes an explanation that is clear, logical, and convincing.}
  \item \textbf{Incorrect}. The map was clear, so tourists could easily find the station. \emph{Clear is more natural for an easy-to-read map; cogent normally describes reasoning or communication intended to persuade.}
  \item \textbf{Correct} \emph{Cogent can appear after found when it describes the reasons.}
  \item \textbf{Incorrect}. He gave a cogent argument for changing the law. \emph{Use the adjective cogent before the noun argument.}
\end{enumerate}

\subsubsection{Collocation practice}
\begin{enumerate}
  \item \textbf{case.} The article makes a cogent case for protecting the wetland. \emph{Make a cogent case for something is an established pattern.}
  \item \textbf{explanation.} The researchers offered a cogent explanation of the unexpected results. \emph{Cogent explanation is a common formal combination.}
  \item \textbf{argument.} The report presents a cogent argument for increasing the sample size. \emph{Cogent argument describes reasoning that is clear, logical, and persuasive.}
\end{enumerate}

\vfill
\begin{center}
\small\color{Muted} Created and compiled by \textbf{Amir Kooshky}\\
\href{https://t.me/KooshkyTOEFL}{Telegram Channel}\quad\textbullet\quad
\href{https://www.instagram.com/kooshkytoefl}{Instagram}\quad\textbullet\quad
\href{https://t.me/KooshkyTOEFL\_pv}{Personal Telegram}
\end{center}
\end{document}
